% --- FONT & TIẾNG VIỆT ---
\usepackage{fontspec}
\setmainfont{Arial} % Chọn font chữ chính (Arial, Times New Roman, Calibri, ...)

% --- TOÁN HỌC & HÌNH ẢNH & BẢNG ---
\usepackage{amsmath}
\usepackage{graphicx}
\usepackage{float}
\usepackage{array}

% --- VẼ HÌNH & KHUNG BÌA ---
\usepackage{tikz}
\usetikzlibrary{calc}

\usepackage{indentfirst} % Ép đoạn văn đầu tiên cũng thụt đầu dòng
\setlength{\parindent}{1cm} % Độ thụt đầu dòng (tùy chỉnh 1cm hoặc 1.27cm chuẩn Word)
\usepackage{setspace}
\onehalfspacing % Dãn dòng 1.5 cho chuẩn báo cáo
\setlength{\parskip}{6pt} % Thêm khoảng cách giữa các đoạn văn (tùy chỉnh từ 6pt - 8pt)

\usepackage{enumitem}
% Cấp 1 (Level 1) - Dấu gạch ngang -
\setlist[itemize,1]{
    label=-,
    itemsep=2pt,
    topsep=4pt,
    leftmargin=*
}

% Cấp 2 (Level 2) - Dấu cộng +
\setlist[itemize,2]{
    label=+,
    itemsep=2pt,
    topsep=2pt,
    leftmargin=1.5em
}

% Package caption để tùy chỉnh tiêu đề hình ảnh và bảng
\usepackage{caption}

% Packges cho bảng nâng cao
\usepackage{array}     % Hỗ trợ định dạng cột nâng cao
\usepackage{booktabs}  % Kẻ đường kẻ bảng chuẩn học thuật (toprule, midrule, bottomrule)
\usepackage{multirow}  % Gộp nhiều dòng trong bảng
\usepackage{makecell}  % Hỗ trợ xuống dòng bên trong 1 ô

% Phải khai báo cuối cùng để tránh xung đột với các package khác
% Dùng để tạo liên kết trong PDF (mục lục, danh sách hình, danh sách bảng)
\usepackage{hyperref} 

\usepackage[table]{xcolor} % Khai báo gói xcolor hỗ trợ tô màu bảng

